\section{Critical Analysis and Team Insights}
\label{sec:analysis}
\subsection{Structural bias, sample size, and computational cost}
CSA assumes that objects share useful mode-specific linear subspaces. Aligned
axes make this plausible for the studied images, but a single Kronecker product
restricts which joint subspaces can be represented. Separate mode matrices do
not imply statistically independent variation. Shorter unfolded fibers give
smaller estimation problems, yet fibers from one image are correlated, not new
independent training images. Compact structure may improve generalization with
scarce data, but the fiber count alone cannot guarantee it.

The clearest benefit is projection storage: CSA needs $\sum_km_kr_k$ entries,
versus $Dd$ for dense PCA, where $D=\prod_km_k$ and $d=\prod_kr_k$.
For Gabor $28\times23\times40\to4\times2\times3$, including an LDA map of the same
$24\to23$ size gives 830 versus 618,792 weights. This is not a peak-memory or universal
runtime claim. Sample-space PCA can use an $N\times N$ eigenproblem
when $N\ll D$ \citep[p.~40, footnote~2]{xu2008csa}; CSA also pays for tensor contractions and
repeated updates.

\subsection{What constitutes convincing evidence}
At equal output dimension and identical centered inputs, unrestricted PCA has
no higher optimal training reconstruction error than its constrained counterpart.
CSA's equal-storage advantage is compatible with this: cheaper shared bases
leave more budget for core coefficients. The budget, not simply the algorithm
name, determines which claim is being tested.

Recognition is a different objective. High-energy lighting variation may be
uninformative about identity, and LDA is not an automatic remedy. The paper's
90.90\% for CSA+LDA is below direct CSA's 94.40\% on ORL; the authors suggest
overfitting \citep[Sec.~IV-C]{xu2008csa}, but their different core dimensions prevent attributing the entire
gap to LDA. Likewise, our metric reversal demonstrates pipeline dependence.
Cross-paper differences are descriptive: the paper's split aggregation is
unspecified, preprocessing and partitions may differ, historical probes
influenced development, and repeated splits reuse images. FERET agreement therefore does not prove a
better implementation or statistically significant superiority.

\subsection{Lessons and targeted improvements}
Our central lesson is to separate three questions: what structure is imposed,
what the optimizer guarantees, and what the evaluation measures. Vectorization
preserves values; CSA's restriction on the projection changes the hypothesis
class. Exact conditional minimization explains stable reconstruction error,
but does not establish a globally best representation or a better classifier.

Future evaluation should freeze preprocessing and metrics, select ranks within
training-only validation, compare both equal-dimension and equal-storage
budgets, and use fresh held-out data. Multiple orthogonal starts can measure
initialization sensitivity, without claiming global optimality. Regularized LDA
or supervised tensor objectives can target discrimination, but our negative
DATER experiment (Appendix~\ref{app:controls}) shows that adding labels alone
is insufficient evidence of improvement. More flexible or nonlinear representations
are possible extensions, with additional capacity and validation costs; they are
not automatic solutions to the Kronecker restriction. Overall, CSA offers a
compact, interpretable structural bias whose recognition benefits remain
conditional on the data and evaluation protocol.
